\exercisesection{交换 Banach 代数}

在以下习题中，所考虑的代数除特别声明外均以 C 为基域。

\begin{enumerate}
\def\labelenumi{\arabic{enumi})}
\item
  设 A 为交换 Banach 代数，I 为 A 的一个极大理想。证明只有两种可能的情形：或者 I 是 A 的某个特征标的核，或者 I 是包含 \(A^{2}\) 的一个超平面。（利用 I，第 154 页的习题 6。）特别地，若 I 不是闭的，则 I 是 A 的包含 \(A^{2}\) 的稠密超平面。（为得到后一情形的一个例子，参见习题 3。）
\item
  设 A 为由趋于 0 的复数序列 \(x = (x_{n})_{n \geqslant 1}\) 组成的 Banach 代数，赋以范数 \(\|x\| = \sup_{i} |x_{i}|\)。设 I 为 A 中具有限支撑的元素所成的集合。则 I 是 A 的一个稠密理想，且不含于任何极大理想之中。（注意 \(A^{2} = A\)，并利用习题 1。）
\item
  设 \(\Lambda\) 为一个集合，\(p \in [1, +\infty[\)，并设 \(\mathbf{A} = \ell^p(\Lambda)\) 为由族 \(x = (x_{\lambda})_{\lambda \in \Lambda}\)（其中 \(x_{\lambda} \in \mathbf{C}\)）组成的集合，且
\end{enumerate}

\[
\| x \| _ {p} = \left(\sum_ {\lambda} | x _ {\lambda} | ^ {p}\right) ^ {1 / p} <   + \infty .
\]

\begin{enumerate}
\def\labelenumi{\alph{enumi})}
\item
  对按分量相加与相乘的运算以及范数 \(x \mapsto \| x \|_p\)，空间 A 是一个 Banach 代数。
\item
  我们有 \(A^{2} \neq A\)，且 \(A^{2}\) 在 A 中稠密。
\item
  A 的非零特征标所成的空间 \(\mathsf{X}(\mathsf{A})\) 典范地等同于赋有离散拓扑的 \(\Lambda\)。（利用 \(\ell^{p}(\Lambda)\) 的对偶是 \(\ell^{q}(\Lambda)\) 这一事实，其中 q 是 p 的共轭指数。）于是 Gelfand 变换成为恒等映射，它在 \(\Lambda\) 上在无穷远处趋于 0 的连续函数空间中的像是稠密的且不是闭的。
\end{enumerate}

\begin{enumerate}
\def\labelenumi{\arabic{enumi})}
\setcounter{enumi}{3}
\tightlist
\item
  设 \((\alpha_{n})_{n\geqslant 0}\) 是 \(\mathbf{R}_{+}^{*}\) 中的序列，满足 \(\alpha_0 = 1\)，对一切整数 \(\alpha_{m + n}\leqslant \alpha_m\alpha_n\) 与 \(m\) 有 \(n\)，且 \(\alpha_{n}^{1 / n}\to 0\)。设 A 为形式级数 \(x = \sum_{n = 0}^{\infty}x_{n}z^{n}\in \mathbf{C}[[z]]\) 中满足下列条件的那些所成的集合：
\end{enumerate}

\[
\| x \| = \sum_ {n = 0} ^ {\infty} | x _ {n} | \alpha_ {n} <   + \infty .
\]

  证明 A 是由 z 拓扑生成的交换幺 Banach 代数，且 A 的唯一极大理想是 \(x \in A\) 中无常数项的元素所成的集合。（注意 z 是拟幂零元。）

\begin{enumerate}
\def\labelenumi{\arabic{enumi})}
\setcounter{enumi}{4}
\item
  在代数 \(\mathbf{C}(\mathrm{X})\)（由 \(\mathbf{C}\) 上的有理分式组成）中，理想 \(\{0\}\) 是极大的，但其余维数不为 1。
\item
  设 \(\Delta\) 为 C 中的闭单位圆盘。证明 \(\mathcal{C}(\Delta)\) 中的可逆元群在 \(\mathcal{C}(\Delta)\) 中不稠密。
\item
  设 A 为交换幺 Banach 代数，\(\chi \in \mathsf{X}(\mathsf{A})\)。对每个 \(x \in \mathsf{A}\)，令 \(\|x\|_{\chi} = |\chi(x)| + \|x - \chi(x)\|\)。证明 \(x \mapsto \|x\|_{\chi}\) 是一个范数，\(\|xy\|_{\chi} \leqslant \|x\|_{\chi}\|y\|_{\chi}\)，\(\|e\|_{\chi} = 1\)，且若 \(\|e\| = 1\)，则 \(\|x\| \leqslant \|x\|_{\chi} \leqslant 3\|x\|\)。
\item
  设 A 为满足 \(\|1\| = 1\) 的交换幺 Banach 代数。设 N 为与给定范数等价、且使 A 仍是 Banach 代数、并使 1 的范数为 1 的范数全体所成的集合。
\end{enumerate}

\begin{enumerate}
\def\labelenumi{\alph{enumi})}
\tightlist
\item
  设 \(x_{0} \in A\) 满足 \(\varrho(x_{0}) < 1\)。对每个 \(x \in A\)，令：
\end{enumerate}

\[
\left\| x \right\| ^ {\prime} = \inf \sum_ {n} \left\| a _ {n} \right\|
\]

其中 \(x\) 取形如 \(a_0 + a_1 x_0 + \cdots + a_n x_0^n\) 的表示（\(a_0, a_1, \ldots, a_n \in \mathrm{A}\)）。证明 \(x \mapsto \|x\|^{\prime}\) 属于 N。（由于 \(\varrho(x_0) < 1\)，存在 \(k\) 使对一切 \(\|x_0^n\| \leqslant k\) 有 \(n\)，从而 \(\|x\| \leqslant k \|x\|^{\prime}\)。）

\begin{enumerate}
\def\labelenumi{\alph{enumi})}
\setcounter{enumi}{1}
\item
  证明 \(\|x_{0}\|^{\prime}\leqslant1\)。
\item
  由 a) 与 b) 推出，对每个 \(x \in A\)，有
\end{enumerate}

\[
\varrho (x) = \inf _ {n \in \mathrm{N}} n (x).
\]

\begin{enumerate}
\def\labelenumi{\alph{enumi})}
\setcounter{enumi}{3}
\tightlist
\item
  若 N 中唯一的元素就是给定的范数，则 \(A = C \cdot 1\)。（利用 c) 与习题 7。）
\end{enumerate}

\begin{enumerate}
\def\labelenumi{\arabic{enumi})}
\setcounter{enumi}{8}
\item
  设 X 为局部紧空间，A 为 Banach 代数 \(\mathcal{C}_{0}(X)\)。设 I 为 A 的一个理想，使得对每个 \(x \in X\)，存在 \(f \in I\) 满足 \(f(x) \neq 0\)。证明 I 包含 X 中具紧支撑的函数所成的理想。
\item
  设 A 为交换 Banach 代数，D 为 A 到 A 中的一个连续导子，即 A 到 A 中的连续线性映射，使得对 \(\mathrm{D}(ab)=(\mathrm{Da})b+a(\mathrm{Db})\) 有 \(a,b\in A\)。
\end{enumerate}

\begin{enumerate}
\def\labelenumi{\alph{enumi})}
\tightlist
\item
  对每个 \(\chi \in \mathsf{X}'(\mathbf{A})\) 与每个 \(\lambda \in \mathbf{C}\)，级数
\end{enumerate}

\[
\varphi_ {\lambda} (a) = \sum_ {n = 0} ^ {\infty} \frac {\lambda^ {n}}{n !} \chi (\mathrm{D} ^ {n} (a))
\]

收敛；\(\lambda \mapsto \varphi_{\lambda}(a)\) 是整函数，而 \(a \mapsto \varphi_{\lambda}(a)\) 是 A 的一个特征标。

\begin{enumerate}
\def\labelenumi{\alph{enumi})}
\setcounter{enumi}{1}
\item
  数 \(\varphi_{\lambda}(a)\) 与 \(\lambda\) 无关。（注意由 \(|\varphi_{\lambda}(a)| \leqslant \|a\|\)) 有 \(a\)。）由此得 \(\chi(\mathrm{D}(a)) = 0\)。
\item
  D 的像含于 A 的根之中。由此给出 I，第 162 页习题 31 的断言 f) 的一个新证明。
\end{enumerate}

\begin{enumerate}
\def\labelenumi{\arabic{enumi})}
\setcounter{enumi}{10}
\tightlist
\item
  设 B 为 {[}0,1{]} 上无限次可微的 C 值函数所成的代数。
\end{enumerate}

\begin{enumerate}
\def\labelenumi{\alph{enumi})}
\tightlist
\item
  设 A 为 B 的一个子代数，赋以一个使 A 成为 Banach 代数的范数。证明存在序列 \((m_{n})_{n\geqslant0}\)（取值于 \(R_{+}\)），使得对一切 \(x\in A\)，有
\end{enumerate}

\[
\sup _ {0 \leqslant t \leqslant 1} | x ^ {(n)} (t) | = \mathrm{O} (m _ {n}).
\]

（设 \(\mathrm{B}_n\) 为 \(n\) 次连续可微函数所成的 Banach 代数，定义在 {[}0,1{]} 上。由 I，第 40 页的命题 9，典范单射 \(\mathrm{A} \to \mathrm{B}_n\) 连续；设 \(\mathrm{M}_n\) 为其范数；若 \(x \in \mathrm{A}\)，则 \(\sup_{0 \leq t \leq 1} |x^{(n)}(t)| \leqslant n!\mathrm{M}_n \| x\|\)。）

\begin{enumerate}
\def\labelenumi{\alph{enumi})}
\setcounter{enumi}{1}
\tightlist
\item
  B 上不存在使它成为 Banach 代数的范数。
\end{enumerate}

\begin{enumerate}
\def\labelenumi{\arabic{enumi})}
\setcounter{enumi}{11}
\tightlist
\item
  设 A 为幺 Banach 代数，\((x_{n})\) 为 A 中趋于元素 \(x \in A\) 的可逆元序列。若序列 \((\varrho(x_{n}^{-1}))\) 有界，且对一切 n 有 \(x_{n}x = xx_{n}\)，则 x 可逆。（由 I，第 38 页命题 8 的推论，\(\varrho(1 - x_{n}^{-1}x) \leqslant \varrho(x_{n}^{-1})\varrho(x_{n} - x)\)，故对充分大的 n，\(x_{n}^{-1}x\) 可逆。）
\end{enumerate}

\begin{enumerate}
\def\labelenumi{¶ \arabic{enumi})}
\setcounter{enumi}{12}
\tightlist
\item
  设 A 为 Banach 代数 \(\ell^{2}(\mathbf{N})\)（习题 3）。设 \(\mathrm{A}_{0}\) 为由具有限支撑的序列组成的子代数。设 B 为 \(\mathrm{A}_{0}\) 与 \(\mathbf{C}\) 的直和，赋以乘法 \((f,\alpha)(g,\beta)=(fg,0)\)（对 \(f,g\in\mathrm{A}_{0},\alpha,\beta\in\mathbf{C}\)）以及范数
\end{enumerate}

\[
\left\| (f, \alpha) \right\| = \sup \left(\left\| f \right\|, \left| \alpha - \sum_ {n} f (n) \right|\right).
\]

\begin{enumerate}
\def\labelenumi{\alph{enumi})}
\item
  完备化代数 \(\widehat{\mathbf{B}}\) 是一个交换 Banach 代数，其根 R 为 \(\mathbf{C} \cdot (0,1)\)，且 \(\widehat{\mathbf{B}}/\mathbf{R}\) 等距同构于 A。
\item
  设 \(\pi\colon\widehat{\mathrm{B}}\to\widehat{\mathrm{B}}/\mathrm{R}\) 为典范态射。不存在 \(\mathrm{B}_{1}\)（\(\widehat{\mathrm{B}}\) 的、与 R 互补的子代数）使 \(\pi|\mathrm{B}_{1}\colon\mathrm{B}_{1}\to\widehat{\mathrm{B}}/\mathrm{R}=\mathrm{A}\) 是同胚。（设 \(\mathrm{B}_{1}\) 是这样一个子代数。设 \(u_{n}\in\mathrm{A}\) 满足 \(u_{n}(n)=1\)、\(u_{n}(m)=0\)（当 \(m\neq n\)）。设 \(e_{n}\in\mathrm{B}_{1}\) 满足 \(\pi(e_{n})=u_{n}\)。证明 \(e_{n}\in\mathrm{A}_{0}\)，且 \(\sum_{n=1}^{\infty}\frac{1}{n}u_{n}\) 在 A 中收敛，而 \(\sum_{n=1}^{\infty}\frac{1}{n}e_{n}\) 在 \(\widehat{\mathrm{B}}\) 中不收敛。）
\end{enumerate}

\begin{enumerate}
\def\labelenumi{\arabic{enumi})}
\setcounter{enumi}{13}
\tightlist
\item
  在 \(\mathbf{C}^2\) 中，设 \(\Omega\) 为由下式定义的紧集：
\end{enumerate}

\[
\frac {1}{2} \leqslant | z _ {1} | ^ {2} + | z _ {2} | ^ {2} \leqslant 1,
\]

并设 A 为 \(\mathcal{C}(\Omega)\) 的由在 \(\mathring{\Omega}\) 中全纯的函数组成的子代数。设 U 为由 \(|z_1|^2 + |z_2|^2 < 1\) 定义的开集。根据 Hartogs 定理（参见 J.-P. DEMAILLY, Complex analytic and differential geometry, Ch. I, th. 3.28），对每个 \(f \in \mathrm{A}\)，存在函数 \(\widetilde{f} \in \mathcal{C}(\overline{\mathrm{U}})\)

，它是唯一的，在 U 中全纯，且在 \(\Omega\) 上与 f 相合。证明 A 到 \(\mathcal{C}(\Omega)\) 中的典范单射是一个等距态射 h，其像是 \(\mathcal{C}(\Omega)\) 的一个全子代数，但 \(\mathsf{X}(h)\) 不是满射的。

\begin{enumerate}
\def\labelenumi{\arabic{enumi})}
\setcounter{enumi}{14}
\item
  设 A 与 B 为交换幺 Banach 代数，\(\varphi\) 为 A 到 B 的保单位元态射，\((x_{\lambda})_{\lambda \in \Lambda}\) 为 A 中一族元素，使得由诸 \(x_{\lambda}\) 生成的 A 的闭全子代数等于 A。为使 \(\mathsf{X}(\varphi)\) 是满射的，充要条件是 \(\mathrm{Sp}_{\mathrm{A}}^{\Lambda}((x_{\lambda})) = \mathrm{Sp}_{\mathrm{B}}^{\Lambda}((\varphi (x_{\lambda})))\)。（利用 I，第 41 页第 6 节的图 (1)，其中右端的箭头由 I，第 43 页的命题 12 知是双射。）
\item
  设 A 为交换 Banach 代数。下列条件等价：
\end{enumerate}

\begin{enumerate}
\def\labelenumi{(\roman{enumi})}
\item
  A 无根，且 \(\mathcal{G}(\mathrm{A})\) 在 \(\mathcal{C}_0(\mathsf{X}(\mathrm{A}))\) 中是闭的；
\item
  \(\mathcal{G}\) 是 A 到 \(\mathcal{G}(\mathrm{A})\) 上的一个同胚；
\item
  存在常数 c \textgreater{} 0，使得对一切 \(\|x\|^{2} \leqslant c\|x^{2}\|\) 有 \(x \in A\)。
\end{enumerate}

\begin{enumerate}
\def\labelenumi{\arabic{enumi})}
\setcounter{enumi}{16}
\item
  设 A 为交换 Banach 代数，I 为 A 的一个闭理想。用 h 表示 I 到 A 中的典范单射。空间 \(\mathsf{X}'(\mathsf{I})\) 典范地等同于 \(\mathsf{X}'(\mathsf{A})\) 关于由 \(\mathsf{X}'(h)\) 所定义的等价关系的商空间。（利用 I，第 153 页习题 4 的 a)。）
\item
  设 A 为交换幺 Banach 代数，\(x\) 为 A 的一个元素。假设由 \(x\) 生成的 A 的闭全子代数等于 A，于是映射 \(\chi \mapsto \chi(x)\) 使 \(\mathsf{X}(\mathsf{A})\) 与 \(\mathrm{Sp}_{\mathsf{A}}(x)\) 得以等同。设 \(\mathcal{H}\) 为函数 \(\mathcal{G}y\)（在 \(\mathrm{Sp}_{\mathsf{A}}(x)\) 上，其中 \(y\) 取遍 A）所成的集合。则 \(\check{\mathrm{S}}_{\mathcal{H}}(\mathrm{Sp}_{\mathsf{A}}(x))\)（INT，IV，§7，第 4 节）是 \(\mathrm{Sp}_{\mathsf{A}}(x)\) 相对于 C 的边界。（为证明 \(\check{\mathrm{S}}_{\mathcal{H}}(\mathrm{Sp}_{\mathsf{A}}(x))\) 含于这个边界，利用最大模原理。反之，设 \(z_0\) 为这个边界上的一点；为证明 \(z_0 \in \check{\mathrm{S}}_{\mathcal{H}}(\mathrm{Sp}_{\mathsf{A}}(x))\)，考虑 \((x - z_1)^{-1}\)，其中 \(z_1\) 充分接近于 \(z_0\)（在 \(\mathbf{C} = \mathrm{Sp}_{\mathsf{A}}(x)\) 中）。）
\item
  设 A 为交换幺 Banach 代数，B 为 A 的一个闭幺子代数，T 为从 X(A) 到 X(B) 中的映射 \(\chi \mapsto \chi|B\)，并设 R 为由 T 在 X(A) 上定义的等价关系。则 T 通过过渡到商，诱导出 X(A)/R 到 T(X(A)) 上的一个同胚，且对每个 \(x \in B\)，函数 \(f = \mathcal{G}_{\mathrm{B}}(x)\) 使 \(|f|\) 在 T(X(A)) 上达到其上确界。
\item
  设 \(\Delta\) 为圆盘 \(|z| \leqslant 1\)（在 \(\mathbf{C}\) 中），A 为由 \(f \in \mathcal{C}(\Delta)\) 中在 \(\Delta\) 的内部全纯的 f 所成的集合，而 \(\mathrm{A}_1\) 为由 A 与函数 \(\mathcal{C}(\Delta)\) 生成的 \(z \mapsto |z|\) 的 Banach 子代数。证明 \(\mathrm{X}(\mathrm{A}_1)\) 同胚于由 \((x_{1}, x_{2}, x_{3}) \in \mathbf{R}^{3}\)（满足 \(x_{1}^{2} + x_{2}^{2} \leqslant x_{3}^{2}\) 与 \(0 \leqslant x_{3} \leqslant 1\)）所成的集合。
\item
  \begin{enumerate}
  \def\labelenumii{\alph{enumii})}
  \tightlist
  \item
    设 \((\mathrm{X}_{\lambda})_{\lambda\in\Lambda}\) 为一族不定元。对每个形式级数 \(f\in\mathbf{C}[[(\mathrm{X}_{\lambda})]]_{\lambda\in\Lambda}\)，用 \(\|f\|\) 表示其系数的绝对值之和。设 A 为使 \(f\in\mathbf{C}[[(\mathrm{X}_{\lambda})]]_{\lambda\in\Lambda}\) 满足 \(\|f\|<+\infty\) 的那些 f 所成的集合。赋以通常的加法与乘法，A 是一个无根的交换幺 Banach 代数。
  \end{enumerate}
\end{enumerate}

\begin{enumerate}
\def\labelenumi{\alph{enumi})}
\setcounter{enumi}{1}
\tightlist
\item
  对每个交换幺 Banach 代数 B，存在一个 a) 中那种类型的代数 A，以及 A 到 B 上的一个连续保单位元态射。
\end{enumerate}

\begin{enumerate}
\def\labelenumi{\arabic{enumi})}
\setcounter{enumi}{21}
\tightlist
\item
  设 \(\Omega\) 为 \(\mathbf{C}^n\) 的一个紧子集。映射
\end{enumerate}

\[
(z _ {1}, \ldots , z _ {n}) \longmapsto (z _ {1}, \ldots , z _ {n}, \overline {{z}} _ {1}, \ldots , \overline {{z}} _ {n})
\]

是从 \(\Omega\) 到 \(\mathbf{C}^{2n}\) 中的一个紧的多项式凸子集上的同胚。（设 \(A = \mathcal{C}(\Omega)\)，\(z_i\) 为 \(\mathbf{C}^n\) 上的坐标函数。则 \((z_i|\Omega, \overline{z}_i|\Omega)_{1 \leqslant i \leqslant n}\) 是 A 的一个拓扑生成组，而所考虑的映射就是由这个拓扑生成组定义的、从 \(X(A) = \Omega\) 到 \(\mathbf{C}^{2n}\) 中的映射。）

\begin{enumerate}
\def\labelenumi{\arabic{enumi})}
\setcounter{enumi}{22}
\tightlist
\item
  用 \(\Omega\) 表示 \((z_1, z_2) \in \mathbf{C}^2\) 中满足 \(|z_1| \leqslant 1\)、\(|z_2| \leqslant 1\) 的元素所成的集合。设 \(\alpha \in [0,1]\)，并用 \(\Omega_\alpha\) 表示 \((z_1, z_2) \in \mathbf{C}^2\) 中满足
\end{enumerate}

\[
| z _ {1} | \leqslant 1, \quad | z _ {2} | \leqslant (1 - \alpha) | z _ {1} | + \alpha .
\]

\begin{enumerate}
\def\labelenumi{\alph{enumi})}
\item
  \(\Omega\) 与 \(\Omega_{\alpha}\) 在 \(\mathbf{C}^2\) 中的余集是连通的。
\item
  \(\Omega\) 是 \(\Omega_{\alpha}\) 的多项式凸包。（若 \((z_1^0, z_2^0) \in \Omega\) 且 \(\mathrm{P} \in \mathbf{C}[z_1, z_2]\)，则存在 \(z \in \mathbf{C}\) 使 \(|z| = 1\) 且 \(|\mathrm{P}(z_1^0, z_2^0)| \leqslant |\mathrm{P}(z, z_2^0)|\)；而 \(|\mathrm{P}(z, z_2^0)| \leqslant \sup_{(z_1, z_2) \in \Omega_{\alpha}} |\mathrm{P}(z_1, z_2)|\)。）
\end{enumerate}

\begin{enumerate}
\def\labelenumi{\arabic{enumi})}
\setcounter{enumi}{23}
\item
  设 \(n \geqslant 1\)。存在 \(\mathbf{C}^{n}\) 的一个闭凸子集，它不是多项式凸的。
\item
  设 A 为一个完备的交换幺局部 m-凸代数（参见 I，第 165 页，习题 37）。
\end{enumerate}

\begin{enumerate}
\def\labelenumi{\alph{enumi})}
\item
  A 的元素 x 可逆，充要条件是 \(\chi(x) \neq 0\) 对 A 的每个连续特征标 \(\chi\) 成立。A 的根是 A 的诸连续特征标的核的交，故是闭的。
\item
  在 I，第 165 页习题 37 的例 e) 中，设 I 为由使对每个充分大的整数 n \(f(n) = 0\)\textgreater{ 0 都有 } 的 f 所成的理想。则每个包含 I 的极大理想都是稠密的且具有无穷余维数。
\end{enumerate}

\begin{enumerate}
\def\labelenumi{\arabic{enumi})}
\setcounter{enumi}{25}
\tightlist
\item
  设 X 为紧拓扑空间，E 为复 Banach 空间，其范数记作 \(\|\cdot\|_{E}\)。用 \(\mathcal{C}(X,E)\) 表示从 X 到 E 中的连续函数所成的复向量空间。
\end{enumerate}

\begin{enumerate}
\def\labelenumi{\alph{enumi})}
\tightlist
\item
  赋以范数
\end{enumerate}

\[
\| f \| = \sup _ {x \in \mathrm{X}} \| f (x) \| _ {\mathrm{E}},
\]

空间 \(\mathcal{C}(\mathrm{X},\mathrm{E})\) 是一个 Banach 空间。

\begin{enumerate}
\def\labelenumi{\alph{enumi})}
\setcounter{enumi}{1}
\tightlist
\item
  设 B 为 \(\mathcal{C}(\mathrm{X},\mathrm{E})\) 的由函数 \(x\mapsto f(x)y\)（其中 \(f\in \mathcal{C}(\mathrm{X})\)、\(y\in \mathrm{E}\)）生成的向量子空间。证明 B 在 \(\mathcal{C}(\mathrm{X},\mathrm{E})\) 中稠密。
\end{enumerate}

在以下内容中，我们考虑一个交换 Banach 代数 A。

\begin{enumerate}
\def\labelenumi{\alph{enumi})}
\setcounter{enumi}{2}
\tightlist
\item
  赋以乘法 \((fg)(x) = f(x)g(x)\)（对一切 \(x\in \mathrm{X}\)），Banach 空间 \(\mathcal{C}(\mathrm{X},\mathrm{A})\) 是一个交换 Banach 代数。
\end{enumerate}

\begin{enumerate}
\def\labelenumi{\alph{enumi})}
\setcounter{enumi}{3}
\tightlist
\item
  每个特征标 \(\chi\in\mathsf{X}(\mathcal{C}(\mathrm{X},\mathrm{A}))\) 都具有形式 \(f\mapsto\widetilde{\chi}(f(x))\)，其中 \(\widetilde{\chi}\in\mathsf{X}(\mathrm{A})\) 且 \(x\in\mathrm{X}\)。（先处理 A 幺且 \(\|1\|=1\) 的情形；考虑 \(\chi\) 在 \(\mathcal{C}(\mathrm{X})\cdot1\) 上的限制，以及它在同态 \(\mathrm{A}\to\mathcal{C}(\mathrm{X},\mathrm{A})\) 的像上的限制，该同态把 \(a\) 映为常值函数 \(x\mapsto a\)。）
\end{enumerate}

\begin{enumerate}
\def\labelenumi{\alph{enumi})}
\setcounter{enumi}{4}
\tightlist
\item
  从 \(\mathsf{X}(\mathrm{A})\times \mathrm{X}\) 到 \(\mathsf{X}(\mathcal{C}(\mathrm{X},\mathrm{A}))\) 中的、把 \((\widetilde{\chi},x)\) 对应到特征标 \(f\mapsto \widetilde{\chi} (f(x))\) 的映射是一个同胚。
\end{enumerate}

\begin{enumerate}
\def\labelenumi{\arabic{enumi})}
\setcounter{enumi}{26}
\tightlist
\item
  设 A 为 C 上的交换幺 Banach 代数。X(A) 的闭子集 F 称为 A 的边界，如果有
\end{enumerate}

\[
\sup _ {\chi \in \mathsf {X} (\mathrm{A})} | \chi (x) | = \sup _ {\chi \in \mathrm{F}} | \chi (x) |
\]

对所有 \(x \in A\) 成立。

\begin{enumerate}
\def\labelenumi{\alph{enumi})}
\item
  将 A 的边界所成的集合按包含关系排序。存在一个极小边界。
\item
  设 S 为 A 的一个极小边界。A 的每个边界都包含 S，因而 S 是 A 的所有边界之交。（设 F 为 A 的一个边界；证明对每个 \(x \in S\)，x 的每个邻域 U 都与 F 相交；为此，利用 S − U 不是 A 的边界这一事实。）
\end{enumerate}

称 S 为 A 的 Shilov 边界，并记作 \(\partial A\)。

\begin{enumerate}
\def\labelenumi{\alph{enumi})}
\setcounter{enumi}{2}
\item
  若 A 的 Gelfand 变换的像在 \(\mathcal{C}(\mathrm{X}(\mathrm{A}))\) 中稠密，则 A 的 Shilov 边界等于 \(\mathrm{X}(\mathrm{A})\)。
\item
  设 A 为单位圆盘 \(\Delta\)（\(z \in C\) 满足 \(|z| \leqslant 1\)）上的、在其内部 \(\Delta\) 全纯的连续函数所成的代数（I，第 20 页的例 9）。则 \(\mathsf{X}(\mathsf{A})\) 典范地等同于 \(\Delta\)（I，第 193 页的习题 6），且 A 的 Shilov 边界等同于单位圆周 U。
\end{enumerate}

\begin{enumerate}
\def\labelenumi{\arabic{enumi})}
\setcounter{enumi}{27}
\item
  设 X 为紧拓扑空间，n 为整数 \(\geqslant 1\)。设 I 为 Banach 代数 \(\mathcal{C}(\mathrm{X},\mathrm{M}_{n}(\mathbf{C}))\) 的一个双边理想。证明存在唯一的闭子集 S ⊂ X，使得 I 恰为 \(f \in \mathcal{C}(\mathrm{X}, \mathrm{M}_{n}(\mathbf{C}))\) 中在 S 上限制为零的 f 所成的集合。
\item
  * 设 X 为紧实流形。给 \(\mathcal{C}^{\infty}(\mathrm{X})\) 赋以函数及其各阶导数一致收敛的拓扑；它是一个 Fréchet 空间。设 A 为 \(\mathcal{C}(\mathrm{X})\) 的一个子代数。假设 A 赋有一个使它成为 Banach 代数的范数 \(f \mapsto \| f\|_{\mathrm{A}}\)，且 A 到 \(\mathcal{C}(\mathrm{X})\) 中的典范单射连续。再假设 \(\mathcal{C}^{\infty}(\mathrm{X})\) 含于 A 中且在 A 中稠密。
\end{enumerate}

\begin{enumerate}
\def\labelenumi{\alph{enumi})}
\tightlist
\item
  存在常数 \(C \geqslant 0\) 与整数 \(k \in N\)，使得对一切 \(f \in \mathcal{C}^{\infty}(X)\)，有
\end{enumerate}

\[
\| f \| _ {\mathrm{A}} \leqslant \operatorname * {C} \sup _ {x \in \mathrm{X}} | f ^ {(k)} (x) |.
\]

\begin{enumerate}
\def\labelenumi{\alph{enumi})}
\setcounter{enumi}{1}
\item
  子代数 A 在 \(\mathcal{C}(\mathrm{X})\) 中是全的。（设 \(f \in \mathrm{A}\) 在 X 上不取零值，并设 \(\delta = \inf_{x \in \mathrm{X}} |f(x)|\)；选取 \(g \in \mathcal{C}^{\infty}(\mathrm{X})\) 使 \(\|f - g\|_{\mathrm{A}} < \delta/3\)，然后在 \(f^{-1} = g^{-1} \sum_{n \in \mathbf{N}} ((g - f)/g)^{n}\) 中写出 \(\mathcal{C}(\mathrm{X})\)，并验证该级数在 A 中收敛。）
\item
  由此给出 Wiener 定理（I，第 38 页的例）的一个新证明。
\end{enumerate}

